\section{Endpoint reconstruction versus process-history reconstruction}

Let $\mathcal W$ be a class of admissible ordered thermodynamic process words and let
\[
 E:\mathcal W\to\operatorname{End}(\Rth)
\]
map each word to its endpoint response operator.

\begin{theorem}[Thermodynamic endpoint-history no-go]
Let $O$ be any observation that factors through the endpoint, $O=\widetilde O\circ E$. If two process words satisfy
\[
 E(w_1)=E(w_2),
\]
then
\[
 O(w_1)=O(w_2).
\]
Therefore no endpoint-only thermodynamic observation can recover all process histories unless the endpoint map is injective on the declared word class.
\end{theorem}

This produces two distinct blindness types.

\begin{description}[leftmargin=2.1em,style=nextline]
\item[Type A: endpoint spectral blindness.] Distinct endpoints share the selected spectral shadow. Marked endpoint observations can repair this blindness.
\item[Type B: erased-history blindness.] Distinct process histories share the same endpoint. Endpoint-only markers cannot repair this blindness; additional history-bearing state is required.
\end{description}

Possible lawful history-bearing augmentations include authenticated prefix responses, internal-position observations, path lifts, or a declared memory transcript. The present paper proves the need for such data but does not construct a universal thermodynamic history representation.

\section{Finite calibration model}

The following exact nilpotent model demonstrates the logical separation between spectral shadow and marked endpoint response. On $\C^3$, set
\[
 X=I+E_{12},
 \qquad
 Y=I+E_{23}.
\]
Then
\[
 XY=I+E_{12}+E_{23}+E_{13},
 \qquad
 YX=I+E_{12}+E_{23},
\]
so
\[
 YX-XY=-E_{13}\ne0.
\]
Nevertheless, \cref{thm:two-factor-blind} gives identical characteristic determinants. With marker $M=E_{31}$,
\[
 \Tr(M(YX-XY))=-1,
\]
and the marked determinant difference appears in the linear coefficient.

This is a finite theorem calibration, not an experimental material model. A physical realization must separately identify the response carrier, generators, metric, marker, admissible parameter range, and target.

\section{Claim boundary and next theorem gates}

The paper proves a finite thermodynamic observation theorem, not a universal constitutive theory. The status is:

\begin{center}
\begin{tabular}{p{0.67\textwidth}p{0.23\textwidth}}
\toprule
Claim & Status\\
\midrule
Typed finite thermodynamic response fibre and metric & Declared datum\\
Two-factor and cyclic characteristic blindness & Proved\\
Cubic connected noncyclic channel & Proved\\
Metric-marked endpoint separation & Proved\\
Complete-bank frame bounds and noisy recovery & Proved\\
Thermodynamic cut-square identity and six-minor expansion & Proved\\
Contractive decoder burden criterion & Proved\\
Degree-five $50/45/5$ collapsed classification & Proved by specialization\\
Minimal five-probe split-occurrence repair & Proved by specialization\\
Target-relative minimum scalar observer count & Proved\\
Endpoint-history no-go & Proved\\
Canonical physical realization of response generators & Open protocol datum\\
Identification of abstract markers with laboratory sensors & Open experimental datum\\
Universal history-faithful thermodynamic representation & Open\\
Infinite-dimensional Fredholm determinant extension & Open\\
Material-independent constitutive law & Not claimed\\
\bottomrule
\end{tabular}
\end{center}

The next development should not merely append more applications. It should prove one of the following missing interfaces:
\begin{enumerate}[label=(\arabic*)]
\item a physically typed construction of ordered response generators from a declared family of thermodynamic interventions;
\item a symmetry-constrained minimal marker bank on the degree-five thermodynamic blind sector;
\item a continuous-time limit with explicit Magnus convergence and error control;
\item a trace-class or regularized determinant theorem for infinite response operators;
\item an experimentally implementable relation between cut-square minors and measured response channels.
\end{enumerate}

\section{Conclusion}

Ordered thermodynamic evolution, endpoint response operators, spectral shadows, marked responses, and process histories are different information layers. Ordinary characteristic data identify some distinct endpoints and erase others. Connected trace channels retain a controlled quotient of noncyclic order memory, while declared markers enlarge the representation and recover blind endpoint residues.

The positive thermodynamic response metric sharpens this recovery theory. The cut-square identity decomposes response energy into the selected scalar channel, unused contractive reserve, and transverse mismatch. For four response channels, six exact minors measure complete metric-normalized coherence. At degree five, block-source collapse creates a five-dimensional spectral defect; five split-occurrence probes repair it exactly and minimally. For an arbitrary thermodynamic target, only the image of the blind kernel under that target must be repaired.

The resulting hierarchy is
\[
\begin{gathered}
 \text{process history}
 \longrightarrow
 \text{endpoint response operator}
 \longrightarrow
 \text{spectral shadow}\\
 \longrightarrow
 \text{metric-marked and cut-square response}.
\end{gathered}
\]
Information can be lost at every arrow. A lawful thermodynamic theory must state which arrow it uses, which blind fibres remain, and which target the observations are intended to preserve.
